\subsection{FRM EXAM 2007---QUESTION 101: evidence of poor 1-day VaR(95\%) estimation}

\textbf{Enunciado}

\emph{A large, international bank has a trading book ... one-day VAR at the 95\% confidence level is USD 50 million.
Which of the following would be the most convincing evidence that the manager is doing a poor job, assuming i.i.d.?}
\begin{enumerate}[label=\alph*.]
\item Over the past 250 days, there are eight exceptions.
\item Over the past 250 days, the largest loss is USD 500 million.
\item Over the past 250 days, the mean loss is USD 60 million.
\item \hl{Over the past 250 days, there is no exception.}
\end{enumerate}

\subsubsection*{Solución}


\textbf{Modelo de excepciones. }Con VaR al 95\%, la probabilidad diaria de excepción es \(p=0.05\).
En \(n=250\) días, el conteo \(N\sim\mathrm{Binomial}(250,0.05)\) y \(\E[N]=np=12.5\).

\textbf{Evidencia más convincente}\\
\begin{itemize}
\item (a) \(N=8\) no es inverosímil (está cerca de 12.5), no es evidencia fuerte de mal modelo.
\item (b) VaR no acota pérdidas máximas; un máximo grande no invalida el VaR.
\item (c) ``mean loss'' no es métrica directa de calibración de VaR.
\item (d) \(N=0\) es extremadamente improbable si el VaR está bien calibrado:
\[
\Prob(N=0)=(1-p)^{250}=0.95^{250}\approx 2.7\times 10^{-6}.
\]
Esto sugiere un VaR mal estimado (probablemente demasiado conservador o mala especificación).
\end{itemize}













